% =============================================================================
% Section 4 — PINN Architecture and Training
% PINN-ONB01 Phase 1 manuscript (IJHMT submission)
% Target length: ~1,080 words (equations, tables and code excluded)
% =============================================================================

\section{PINN Architecture and Training}\label{sec:architecture}

The conditioned PINN couples a surface-descriptor encoder, a
one-dimensional conduction backbone, three task-specific heads, and a
five-term composite loss. The loss imposes the priors of
\cref{sec:formulation} as soft constraints. The total parameter count
is \num{24005}. The count was deliberately constrained because only \num{1361}
boiling-curve points were available, decreasing to \num{1167} after
the FC-77 holdout. Hyperparameters for
the final \texttt{baseline\_phaseDbal} model are summarized in
\cref{tab:hyperparameters}.

% -----------------------------------------------------------------------------
\subsection{Surface encoder and PINN backbone}\label{subsec:encoder_backbone}\label{subsec:encoder}\label{subsec:backbone}
% -----------------------------------------------------------------------------

The encoder maps the partially observed surface descriptor of each
heater to a fixed latent $z_s \in \mathbb{R}^{8}$. Its input
concatenates nine numeric channels with a four-dimensional embedding
of the paper category. The numeric channels are
(0) $\log_{10}(R_a/\SI{1}{\micro\meter})$,
(1) $\log_{10}(R_a/L_c)$,
(2) $\cos\theta$,
(3) $\theta/\pi$,
(4) a binary $\theta$-availability mask,
(5) $\log_{10}(r_c/L_c)$,
(6) an availability mask for $r_c$,
(7) $\log_{10}(N_s)/6$ with $N_s$ in \si{\per\centi\meter\squared}, and
(8) an availability mask for $N_s$.
The two complementary roughness channels span four decades while
preserving the capillary scaling of \cref{subsec:nondim}. Channels
(2)--(4) carry the contact-angle dependence of Young's equation and
the Hsu criterion. The masks distinguish ``not measured'' from ``measured zero''. A categorical embedding with a vocabulary of ten
entries, indexed by source paper and heater material, absorbs
paper-level biases without coupling to fluid identity. The full
category list is provided in the supplementary material. Scalar layer
normalization renormalizes only over observed entries (using the
masks) and then feeds a three-layer MLP of width \num{64} with GELU
activations and dropout \num{0.1}. The output is projected to
$\mathbb{R}^{8}$ and bounded by $\tanh$. The encoder contains
\num{6154} trainable parameters.

The backbone is a fully connected network with five hidden layers of
width \num{64} and $\tanh$ activations; the smooth activation keeps
the second-order derivatives required by \cref{eq:laplace} well-defined
under automatic differentiation. The backbone receives
$z^{*} \in [0,L^{*}]$, the surface latent $z_s$, and the operating
descriptor $(q^{*}_{\mathrm{appl}}, \Delta T^{*}_{\mathrm{sub}})$. We
implemented two conditioning strategies. The \texttt{concat} variant
stacks all inputs into an \num{11}-dimensional vector; the
\texttt{film} variant --- adopted in the final model --- injects $z_s$
at every hidden layer through a Feature-wise Linear
Modulation~\citep{perez2018film} affine transform
$h \leftarrow \gamma(z_s)\odot h + \beta(z_s)$. FiLM couples the
surface latent multiplicatively while using fewer parameters than the
concatenation variant. Three two-layer task heads, each with hidden
width \num{32}, branch from the final hidden representation: one head
predicts $T^{*}$ for the PDE and boundary residuals, while the other
two predict $\Delta T^{*}_{\mathrm{ONB}}$ and $q^{*}_{\mathrm{ONB}}$
through softplus activations. The backbone, FiLM generators, and heads
contribute \num{17851} parameters. These components yield a total of \num{24005} parameters. The architecture follows the
\citet{raissi2019} template generalized through surface conditioning.

% -----------------------------------------------------------------------------
\subsection{Composite loss and training pipeline}\label{subsec:loss_training}\label{subsec:loss}\label{subsec:phases}\label{subsec:impl}
% -----------------------------------------------------------------------------

The composite objective comprises PDE, boundary, data, nucleation-prior,
and monotonicity-prior terms, each expressed in the nondimensional
variables of \cref{subsec:nondim}.

The PDE residual enforces \cref{eq:laplace} through automatic
differentiation on $N_{\mathrm{c}}=\num{2000}$ LHS collocation
points~\citep{mckay1979lhs} that are refreshed every \num{100} epochs:
%
\begin{equation}\label{eq:L_cond}
\mathcal{L}_{\mathrm{cond}}
\;=\;
\frac{1}{N_{\mathrm{c}}}
\sum_{i=1}^{N_{\mathrm{c}}}
\left|
\frac{\partial^{2} T^{*}}{\partial {z^{*}}^{2}}\bigg|_{z^{*}_{i}}
\right|^{2},
\end{equation}
%
The boundary loss combines the Neumann condition at $z^{*}=0$ with the
Dirichlet-type closure at $z^{*}=L^{*}$,
%
\begin{equation}\label{eq:L_BC}
\mathcal{L}_{\mathrm{BC}}
\;=\;
\frac{1}{M}\sum_{i=1}^{M}
\bigl(-k^{*}\,\partial_{z^{*}} T^{*}\big|_{0}
       - q^{*}_{\mathrm{appl},i}\bigr)^{2}
\;+\;
\frac{1}{M}\sum_{i=1}^{M}
\bigl(T^{*}(L^{*}) - T^{*}_{\mathrm{bulk},i}\bigr)^{2}.
\end{equation}
%
The supervised data term penalizes prediction errors in both ONB
superheat and heat flux across the \num{1361} curated points:
%
\begin{equation}\label{eq:L_data}
\mathcal{L}_{\mathrm{data}}
\;=\;
\frac{w_{\Delta T}}{N}
\sum_{n=1}^{N}
\bigl(
\Delta T^{*,\mathrm{pred}}_{\mathrm{ONB},n} - \Delta T^{*,\mathrm{obs}}_{\mathrm{wall},n}
\bigr)^{2}
+
\frac{w_{q}}{N}
\sum_{n=1}^{N}
\bigl(
q^{*,\mathrm{pred}}_{\mathrm{ONB},n} - q^{*,\mathrm{obs}}_{n}
\bigr)^{2},
\end{equation}
%
with $w_{\Delta T}$ and $w_{q}$ chosen so that the contributions are
comparable at initialization. The Hsu
term~\citep{hsu1962,basu2002} converts \cref{eq:hsu_discriminant} into
a one-sided hinge. It penalizes predictions for which the discriminant
$D_{\mathrm{Hsu}}(\Delta T^{\mathrm{pred}}_{\mathrm{wall}}, q^{\mathrm{pred}})$
falls below a margin $\varepsilon$:
%
\begin{equation}\label{eq:L_Hsu}
\mathcal{L}_{\mathrm{Hsu}}
\;=\;
\frac{1}{N}
\sum_{n=1}^{N}
\bigl[
\max\bigl(0,\;\varepsilon - D_{\mathrm{Hsu}}(\Delta T^{\mathrm{pred}}_{\mathrm{wall},n},
                                              q^{\mathrm{pred}}_{n})\bigr)
\bigr]^{2}.
\end{equation}
%
We set $\varepsilon=\num{0.05}$. The weight $w_{\mathrm{ONB}}$ was held
at $\num{0.10}$ throughout fine-tuning. The implementation supports a
linear ramp over the first \num{2000} Adam epochs; Bayesian
optimization, however, converged on start and end values that differ
by less than \SI{2}{\percent} ($\num{0.1000}$ and $\num{0.1014}$),
rendering the schedule effectively constant.

The monotonicity term penalizes violations of two established trends.
$\Delta T_{\mathrm{ONB}}$ should decrease with $R_a$ and with $\theta$.
We query the encoder on a counterfactual sweep of $B=\num{20}$ points.
The sweep is log-uniform on
$R_a \in [\SI{0.01}{\micro\meter},\SI{50}{\micro\meter}]$ and linear on
$\theta \in [\SI{10}{\degree},\SI{170}{\degree}]$. Other descriptors
are held at batch medians:
%
\begin{equation}\label{eq:L_mono}
\mathcal{L}_{\mathrm{mono},x}
\;=\;
\frac{1}{B-1}
\sum_{b=1}^{B-1}
\bigl[
\max\bigl(0,\;
\Delta T_{\mathrm{ONB}}(x_{b+1}) - \Delta T_{\mathrm{ONB}}(x_{b})
\bigr)
\bigr]^{2},
\qquad x \in \{R_a, \theta\}.
\end{equation}
%
The counterfactual sweeps regularize unobserved directions without
contradicting supervised data.

The full objective combines the five terms with fixed scalar weights,
%
\begin{equation}\label{eq:L_total}
\begin{aligned}
\mathcal{L} = {}
& w_{\mathrm{cond}}\,\mathcal{L}_{\mathrm{cond}}
+ w_{\mathrm{BC}}\bigl(\mathcal{L}_{\mathrm{BC},0} + \mathcal{L}_{\mathrm{BC},L}\bigr)
+ w_{\mathrm{data}}\bigl(\mathcal{L}_{\Delta T} + \mathcal{L}_{q}\bigr) \\
& {} + w_{\mathrm{ONB}}\,\mathcal{L}_{\mathrm{Hsu}}
+ w_{\mathrm{mono},R_a}\,\mathcal{L}_{\mathrm{mono},R_a}
+ w_{\mathrm{mono},\theta}\,\mathcal{L}_{\mathrm{mono},\theta}.
\end{aligned}
\end{equation}
%
The optimized weights are $w_{\mathrm{cond}}=\num{2.00}$,
$w_{\mathrm{BC}}=\num{3.89}$, $w_{\mathrm{data}}=\num{35.57}$,
$w_{\mathrm{ONB}}=\num{0.10}$, $w_{\mathrm{mono},R_a}=\num{2.50}$, and
$w_{\mathrm{mono},\theta}=\num{1.00}$ (\cref{tab:hyperparameters}). The
balance is data-dominant, reflecting ONB sparsity. The Hsu hinge acts
as soft consistency rather than a fitting target.

We trained the model in three sequential stages, each activating a
richer subset of the loss than the preceding one. The schedule proved
more stable than optimizing \cref{eq:L_total} jointly. The \emph{analytical warm-up}
ran for \num{500} Adam epochs at $\eta=\num{1e-3}$~\citep{kingma2014adam}
using only PDE and boundary terms with $z_s=0$. The conduction
residual converged below \num{e-5} in $\sim$\SI{7}{\second} on a single
Apple~M1 CPU. The \emph{experimental fine-tuning} stage then activated
the full loss. It ran \num{2000} Adam epochs at $\eta=\num{2.6e-3}$
with constant Hsu weight $w_{\mathrm{ONB}}=\num{0.10}$, gradient
clipping, and non-finite-step rollback. Up to \num{1000} L-BFGS iterations with
strong Wolfe line search~\citep{liu1989lbfgs} followed, applied to the
backbone only with the heads frozen. This stage completed in
$\sim$\SIrange{17}{20}{\minute}. The \emph{inverse stage} freezes the
trained network. It introduces a per-surface $r_c$ as a learnable
parameter, which we optimized over \num{500} Adam iterations at
$\eta=\num{1e-3}$ to maximize the Hsu-likelihood of observed ONB
points. We selected hyperparameters via an Optuna TPE search over
\num{20} trials. The best configuration was rebalanced to give the
final model used throughout \cref{sec:results}.

We implemented the network in PyTorch~2 in single
precision~\citep{paszke2019pytorch}, with a custom training loop to
support per-phase parameter freezing and rollback. Mini-batches
of size \num{32} yield \num{36} batches per epoch on the
\num{1167}-point training pool. We split the dataset at the surface
level (\SI{80}{\percent} train, \SI{10}{\percent} validation,
\SI{10}{\percent} test); because no surface appears in more than one
split, the protocol prevents data leakage. Collocation points are redrawn every \num{100}
epochs (LHS seed $31337\cdot\text{epoch}$). Parameter initialization
and split seeds are fixed at \num{42}. Early stopping (patience
\num{200}) rarely triggers. All experiments ran on an Apple~M1 host
with eight CPU cores and \SI{16}{\giga\byte} RAM. Runs, trials, and
checkpoints are logged via MLflow. We release the code, dataset, and
the final \texttt{baseline\_phaseDbal} checkpoint so the pipeline can
be re-executed from \cref{tab:hyperparameters}.

% -----------------------------------------------------------------------------
\begin{table}[!htbp]
\caption{Hyperparameters and weight settings for the final
  \texttt{baseline\_phaseDbal} model.}
\label{tab:hyperparameters}
\small
\centering
\begin{tabular}{lll}
\hline
Category & Parameter & Value \\
\hline
\multirow{4}{*}{Architecture}
  & Conditioning      & FiLM \\
  & Latent dim        & 8 \\
  & Hidden dim        & 64 \\
  & Layers (backbone) & 5 \\
\hline
\multirow{2}{*}{Total params}
  & Encoder           & \num{6154} \\
  & Backbone + heads  & \num{17851} \\
\hline
\multirow{6}{*}{Loss weights}
  & $w_{\mathrm{cond}}$        & \num{2.00} \\
  & $w_{\mathrm{BC}}$          & \num{3.89} \\
  & $w_{\mathrm{data}}$        & \num{35.57} \\
  & $w_{\mathrm{ONB}}$ (Hsu)   & \num{0.10} \\
  & $w_{\mathrm{mono},R_a}$    & \num{2.50} \\
  & $w_{\mathrm{mono},\theta}$ & \num{1.00} \\
\hline
\multirow{4}{*}{Training}
  & Warm-up Adam epochs    & \num{500} \\
  & Fine-tune Adam epochs  & \num{2000} \\
  & Fine-tune L-BFGS iter. & \num{1000} \\
  & $\eta_{\mathrm{Adam}}$ & \num{2.6e-3} \\
\hline
\end{tabular}
\end{table}
% -----------------------------------------------------------------------------
